@@M0@@ that is representable by an affine scheme over @@M1@@. Now @@M2@@ is defined by an invertible sheaf @@M3@@ on @@M4@@ equipped with a rigidification @@M5@@ for @@M6@@. Likewise, a point @@M7@@ corresponds to an invertible sheaf @@M8@@ on @@M9@@ equipped with a rigidification @@M10@@. The image of @@M11@@ in @@M12@@ is @@M13@@ if and only if there is an isomorphism of rigidified sheaves @@M14@@ where @@M15@@ is the inclusion for @@M16@@. Observe that these isomorphisms define an isomorphism @@M17@@ such that @@M18@@.

Conversely, given an isomorphism @@M19@@, @@M20@@, such that @@M21@@, there is a rigidified invertible sheaf @@M22@@ on @@M23@@ obtained by gluing along @@M24@@; explicitly, @@M25@@ is the kernel of the canonical homomorphism @@M26@@ defined by @@M27@@. This @@M28@@ then corresponds to a point @@M29@@ of @@M30@@ mapping to @@M31@@.

Finally, let @@M32@@ be the functor defined, as @@M33@@ runs through the schemes over @@M34@@, by the sets of isomorphisms @@M35@@ such that @@M36@@. We have just established an isomorphism between the inverse image of @@M37@@ in @@M38@@ and the functor @@M39@@. It therefore suffices to apply the following lemma (taking @@M40@@):

\sgaStableTarget{sga6.e12.lem.2-4}\textbf{Lemma 2.4.} \uline{Let} @@M41@@ \uline{be a flat, proper scheme of finite presentation over} @@M42@@\uline{. Let} @@M43@@ \uline{be two invertible sheaves on} @@M44@@, \uline{and} @@M45@@ \uline{the functor of isomorphisms} @@M46@@ \uline{(for} @@M47@@\uline{). Then}
\begin{enumerate}[label=\roman*)]
\item @@M48@@ \uline{is representable by an affine scheme over} @@M49@@ \uline{of finite presentation.}
\item \uline{If, moreover,} @@M50@@ \uline{and} @@M51@@ \uline{are equipped with rigidifications (relative to a}

